\documentclass[11pt]{article}

\usepackage[a4paper,margin=1in]{geometry}
\usepackage{amsmath}
\usepackage{amssymb}
\usepackage{amsthm}
\usepackage{mathtools}
\usepackage{listings}
\usepackage{xcolor}
\usepackage{graphicx}
\usepackage{booktabs}
\usepackage[hidelinks]{hyperref}

% -- Code listing style for Python snippets --
\definecolor{codegray}{gray}{0.95}
\definecolor{codegreen}{rgb}{0.0,0.5,0.0}
\definecolor{codepurple}{rgb}{0.58,0,0.82}
\lstset{
  basicstyle=\ttfamily\small,
  backgroundcolor=\color{codegray},
  keywordstyle=\color{blue},
  commentstyle=\color{codegreen},
  stringstyle=\color{codepurple},
  numbers=left,
  numberstyle=\tiny,
  breaklines=true,
  frame=single,
  language=Python,
  showstringspaces=false,
}

\title{Implementing an AIXItl Agent in a Toy Environment:\\ A Technical Report}
\author{Ivan Yaretskyi}
\date{\today}

\begin{document}
\maketitle

\begin{abstract}
This report documents the design and implementation of an AIXItl agent -- a
computable approximation of Hutter's universal artificial intelligence agent AIXI --
together with the machinery required to build it: Turing and universal Turing machines,
a bounded Kolmogorov complexity calculator, a Solomonoff universal predictor with a
stratified Monte Carlo sampler over its hypothesis space, and (in future work) an
expectimax planner operating over a toy environment. The work is organised as a sequence
of self-contained stages, each described in its own section, together with a retrospective
on the development process itself.
\end{abstract}

\section{Introduction and Motivation}
\label{sec:introduction}

AIXI is a universal artificial intelligence agent proposed by Marcus Hutter. At its core,
it extends Solomonoff's theory of universal induction -- which provides an optimal method
for predicting the next output of any unknown computable process -- to the case of an
active agent that takes actions and receives rewards. Classical reinforcement learning
algorithms rely on structural assumptions about the environment, such as the Markov
property, stationarity, or a known state space. AIXI requires none of these. It instead
maintains a Bayesian mixture over all computable environments, weighted by their
Kolmogorov complexity, and selects actions by planning with an expectimax tree over that
mixture. It is therefore both parameter-free and assumption-free: no domain-specific
tuning is required.

The corresponding drawback is that AIXI is uncomputable. Evaluating the universal mixture
requires simulating infinitely many programs, including non-halting ones, which is
equivalent to solving the halting problem. To circumvent this, I implement \textbf{AIXItl},
a tractable approximation that restricts the hypothesis space to programs of length at
most $l$, each simulated for at most $t$ steps. This preserves the theoretical guarantees
of AIXI for any environment whose generating program fits within these bounds, while
remaining implementable on ordinary hardware. Alternative routes to tractability
exist -- most notably MC-AIXI, which employs Monte Carlo sampling within the Solomonoff
predictor -- and, as later sections show, the practical demands of running this project on
consumer hardware pushed the implementation in a similar direction independently.

\section{Turing Machines and the Universal Turing Machine}
\label{sec:turing}

\subsection{Turing Machines}
\label{sec:tm}

A Turing machine is a fundamental model of computation. It is Turing-complete --
capable of expressing any algorithm that can be expressed at all -- which makes it the
natural substrate for the induction process described in later sections.

The implementation in this project uses a 6-tuple definition,
\[
  (Q,\ \Gamma,\ \delta,\ q_{\text{start}},\ q_{\text{halt}},\ b),
\]
where $Q$ is the finite set of states; $\Gamma$ is the tape alphabet, with a designated
blank symbol at index $b$; $\delta$ is the transition function
\[
  \delta : Q \times \Gamma \to Q \times \Gamma \times \{L, R\} \times \{0,1\},
\]
mapping a (state, tape symbol) pair to a new state, a symbol to write, a direction to move
the head, and an \emph{emit} flag discussed in Section~\ref{sec:emission}; and
$q_{\text{start}}, q_{\text{halt}} \in Q$ are the start and halt states. Computation is
deterministic and stops the moment the machine transitions into $q_{\text{halt}}$.

This departs from the standard 7-tuple definition
$(Q, \Sigma, \Gamma, \delta, q_{\text{start}}, q_{\text{accept}}, q_{\text{reject}})$ found
in the theory-of-computation literature, which distinguishes an accepting halt state from a
rejecting one. Section~\ref{sec:halt-merge} explains why this distinction was dropped, and
in what precise sense the two models remain equivalent for the purposes of this project.

\subsection{Why a Single Halt State}
\label{sec:halt-merge}

The standard 7-tuple model requires $q_{\text{accept}} \neq q_{\text{reject}}$, since the
two states carry distinct meaning. Requiring two distinct halting states has a specific,
damaging consequence for a machine with only $n=2$ states: \emph{both} available states
must be halting states, leaving no state a transition can lead to other than immediate
termination. A 2-state machine under the 7-tuple model therefore always halts after reading
a single symbol, with no capacity to loop, branch on more than one symbol, or compute
anything beyond a single-step response.

This matters because 2-state machines are also the cheapest to describe: their encoded
length is the shortest of any class (Section~\ref{sec:sampling} shows this class alone
accounts for over half of all prior probability mass once the emit-bit scheme of
Section~\ref{sec:emission} is in place). Under the 7-tuple model, the entire class
contributing the most prior mass is guaranteed to be computationally inert -- a large
fraction of the hypothesis space that the Solomonoff prior favours most heavily is, by
construction, incapable of representing anything interesting.

Collapsing $q_{\text{accept}}$ and $q_{\text{reject}}$ into one $q_{\text{halt}}$ frees up
one of the two states in the $n=2$ case: one state can now serve purely as a working state,
able to transition to itself and process more than one input symbol before halting.
Structurally, a single halt state is a strict special case of the general model, not a
different one: any 6-tuple machine embeds into a standard 7-tuple machine by declaring
$q_{\text{halt}}$ the accept state and adding one additional, permanently unreachable
reject state. The embedding is not reversible -- the 6-tuple model has no way to represent
an \emph{explicit} rejection as distinct from an unbounded computation -- but this
asymmetry costs nothing here. The predictor described in Section~\ref{sec:predictor} never
queries whether a candidate program accepted or rejected; it only cares about the sequence
of bits a program emits while running. Collapsing accept and reject into one halt state
therefore loses no information the rest of the system uses, while directly repairing the
most damaging structural defect in the smallest, most heavily weighted machine class.

\subsection{Emitting Output}
\label{sec:emission}

An earlier version of this project treated the machine's tape, read left to right after
halting (or after the step budget was exhausted), as its output. Two problems followed
from this choice.

First, a two-way infinite tape allows the head to move left of its starting position.
Measurement on a large batch of randomly generated machines showed that roughly one in
fifty programs, when compared at two different step budgets, did not have their shorter-run
output as a prefix of their longer-run output: intervening leftward head movement had
overwritten a cell that had already been read as part of the "output" at the earlier
checkpoint. This silently violates the property the predictor depends on most -- that a
program's prediction of bit $n$, once made, does not later change.

Second, and more fundamentally, restricting attention to \emph{halted} programs discards
every machine still computing when the step budget is reached. Measurement showed that
halting machines average only one to two bits of usable output before terminating, while
machines still running at the step limit average forty to a hundred. Since the predictor's
whole purpose is to find programs that keep agreeing with an ongoing, possibly unbounded
environment, discarding every program that has not yet halted discards almost every program
capable of saying anything useful.

The fix was to give the transition function an explicit \emph{emit} flag. When a transition
fires with this flag set, the machine appends the symbol it just wrote (reduced modulo~2, so
the emitted stream is binary regardless of the tape alphabet's size) to a separate,
append-only output list, entirely distinct from the tape. This list is never revised: once a
bit has been emitted, no later transition can alter it, which eliminates the prefix-violation
problem by construction rather than by patching around it after the fact. It also means a
program's status at the step budget -- halted or still running -- is now irrelevant to
whether its output so far is usable: a non-halting machine has still emitted a perfectly
valid, permanently fixed prefix, and is exactly as legitimate a hypothesis as one that has
halted.

This construction is a restricted instance of the standard equivalence between a
single-tape and a multi-tape Turing machine: one ordinary read/write work tape, and one
write-only, append-only output tape whose head only ever advances. The general
multi-tape-to-single-tape reduction (interleaving each tape's cells on tracks of one
physical tape) is more machinery than is needed here, precisely because the output tape's
head never moves backward: its entire history can be represented as a simple growing list
rather than requiring the two-way-infinite bookkeeping a genuine second work tape would
need. The implementation reflects this directly -- \texttt{self.output} is a plain Python
list, appended to and never otherwise modified.

\subsection{The Universal Turing Machine}
\label{sec:utm}

A Universal Turing Machine (UTM) is a Turing machine capable of simulating any other
Turing machine. It takes as input a binary encoding of some machine $M$ together with an
input string $w$, and produces the same behaviour that $M$ would exhibit on $w$. Crucially,
$M$ does not run as a separate process -- it exists entirely as data on the UTM's tape.
The UTM reads $M$'s transition table one rule at a time, determines the applicable
transition for the current simulated state and symbol, and applies it, updating a region
of its own tape that represents $M$'s tape (and, per Section~\ref{sec:emission}, a separate
list representing $M$'s output).

For a UTM to simulate an arbitrary machine $M$, it must be possible to express $M$ as a
finite binary string. The convention for doing so is called an \emph{encoding scheme}. The
choice of encoding scheme is a design decision: any scheme that is unambiguous and
complete is valid, since the Invariance Theorem (Section~\ref{sec:invariance}) guarantees
that Kolmogorov complexity changes by at most an additive constant between any two
reasonable schemes.

\subsection{Encoding Scheme}
\label{sec:encoding}

The encoding scheme used in this project represents a Turing machine as a self-delimiting
binary string structured as follows.

\paragraph{Header.} The string opens with two unary fields. The number of states $n$ is
encoded as $n$ ones followed by a zero; the size of the tape alphabet $k$ is encoded the
same way. These fields are self-delimiting -- a decoder reads ones until it sees a zero,
requiring no external length information.

\paragraph{Field widths.} Once $n$ and $k$ are known, all subsequent binary fields have
fixed, computable widths:
\begin{align*}
  \texttt{state\_bits}  &= \lceil \log_2 n \rceil && \text{(minimum bits to index any state)},\\
  \texttt{symbol\_bits} &= \lceil \log_2 k \rceil && \text{(minimum bits to index any tape symbol)}.
\end{align*}

\paragraph{Special states.} The indices of $q_{\text{start}}$ and $q_{\text{halt}}$ follow
immediately, each encoded in \texttt{state\_bits} bits. This is one field shorter than the
7-tuple scheme, which required three special-state indices rather than two.

\paragraph{Transition table.} The transition function $\delta$ is stored as a flat table
of $n \times k$ entries in lexicographic order of (current state, current symbol) pairs.
The input pair is implicit from the entry's position; only the output is stored. Each
entry consists of four fields:
\begin{center}
\begin{tabular}{ll}
  \toprule
  Field & Width \\
  \midrule
  \texttt{new\_state}  & \texttt{state\_bits} bits \\
  \texttt{new\_symbol} & \texttt{symbol\_bits} bits \\
  \texttt{direction}   & 1 bit ($0 = $ Left, $1 = $ Right) \\
  \texttt{emit}         & 1 bit ($0 = $ silent, $1 = $ append to output) \\
  \bottomrule
\end{tabular}
\end{center}
No delimiters are required. A decoder reads exactly
$(\texttt{state\_bits} + \texttt{symbol\_bits} + 2)$ bits per entry and reconstructs each
transition rule from the entry's position in the table. This is one bit longer per entry
than the 7-tuple scheme, which had no emit flag.

\paragraph{Blank symbol.} Finally, the index of the blank symbol is appended in
\texttt{symbol\_bits} bits.

The net effect of dropping one special-state field while adding one bit to every
transition entry is, for small machines, a slightly \emph{longer} minimum encoding than
the old scheme (25 bits for the smallest machine under this scheme, versus 22 bits under
the old one). This is not the point of the change: the point, developed fully in
Section~\ref{sec:sampling}, is that the smallest class becomes \emph{useful} rather than
merely \emph{short} -- a distinction that matters far more once machines are weighted by
both their prior probability and their capacity to say anything about an unfolding
sequence.

\subsection{Worked Example}
\label{sec:encoding-example}

Consider the machine used as the canonical worked example throughout this project's test
suite: a 3-state, 2-symbol machine that emits the alternating sequence
$0,1,0,1,0,1,\ldots$ forever.
\begin{center}
\begin{tabular}{ll}
  \toprule
  States    & $Q = \{0, 1, 2\}$ \\
  Alphabet  & $k = 2$ symbols \\
  Blank     & index $0$ \\
  Start     & state $0$ \\
  Halt      & state $2$ (unreachable in normal operation) \\
  \bottomrule
\end{tabular}
\end{center}
States $0$ and $1$ ping-pong forever, each writing a fixed symbol to the tape, emitting it,
and handing control to the other, regardless of what is read:
\[
  \delta(0,\cdot) = (1,\,0,\,R,\,\text{emit}), \qquad
  \delta(1,\cdot) = (0,\,1,\,R,\,\text{emit}).
\]
State $2$ is never reached in practice; its two transition-table rows are included only
because the encoding requires an entry for every (state, symbol) pair, and are filled with
an arbitrary silent self-loop.

\paragraph{Computing field widths.}
\begin{align*}
  \texttt{state\_bits}  &= \lceil \log_2 3 \rceil = 2,\\
  \texttt{symbol\_bits} &= \lceil \log_2 2 \rceil = 1.
\end{align*}

\paragraph{Header.}
\begin{center}
\begin{tabular}{ll}
  States (3)   & \texttt{111 0} \\
  Symbols (2)  & \texttt{11 0} \\
\end{tabular}
\end{center}

\paragraph{Special states.}
\begin{center}
\begin{tabular}{ll}
  \texttt{start = 0}  & \texttt{00} \\
  \texttt{halt = 2}   & \texttt{10} \\
\end{tabular}
\end{center}

\paragraph{Transition table} (6 entries, one per (state, symbol) pair in order):
\begin{center}
\begin{tabular}{lll}
  \toprule
  Input pair & Output $(s', \sigma', d, e)$ & Encoding \\
  \midrule
  (state 0, symbol 0) & $(1, 0, R, \text{emit})$ & \texttt{01 0 1 1} \\
  (state 0, symbol 1) & $(1, 0, R, \text{emit})$ & \texttt{01 0 1 1} \\
  (state 1, symbol 0) & $(0, 1, R, \text{emit})$ & \texttt{00 1 1 1} \\
  (state 1, symbol 1) & $(0, 1, R, \text{emit})$ & \texttt{00 1 1 1} \\
  (state 2, symbol 0) & \multicolumn{2}{l}{unreachable halt-state row;} \\
                       & \multicolumn{2}{l}{filled with $(2,0,R,\text{silent})$: \texttt{10 0 1 0}} \\
  (state 2, symbol 1) & \multicolumn{2}{l}{same filler: \texttt{10 0 1 0}} \\
  \bottomrule
\end{tabular}
\end{center}

\paragraph{Blank index.}
\begin{center}
\begin{tabular}{ll}
  \texttt{blank = 0} & \texttt{0} \\
\end{tabular}
\end{center}

\paragraph{Concatenated encoding.}
\begin{verbatim}
1110 110 00 10
01011 01011 00111 00111 10010 10010
0
\end{verbatim}
Removing spaces gives the final 42-bit binary string:
\begin{verbatim}
111011000100101101011001110011110010100100
\end{verbatim}
This encoding, and the machine's behaviour, were verified directly: decoding it and
running it for ten steps reproduces the transition table above and emits exactly
\texttt{0101010101}, confirming both the encoder and decoder against each other and
against the intended machine.

\section{Kolmogorov Complexity}
\label{sec:kolmogorov}

\subsection{Definition and Incomputability}

The (prefix) Kolmogorov complexity of a finite binary string $x$, relative to a fixed
universal machine $U$, is the length of the shortest program that makes $U$ halt having
emitted exactly $x$:
\[
  K(x) = \min\{\, |p| : U(p) \text{ halts with output } x \,\}.
\]
$K$ is not computable. If it were, one could construct a program that enumerates
candidate strings and, for each, computes its complexity, halting on the first string
whose complexity exceeds some threshold $n$. Such a program would itself have a fixed,
short description -- of order $\log n$ bits -- yet by construction it outputs a string
whose complexity is greater than $n$, a contradiction once $n$ exceeds the program's own
length. The same fact can be seen directly from the halting problem: identifying the
\emph{shortest} halting program for $x$ requires being able to rule out every shorter
candidate, and ruling out a candidate that never halts is exactly the halting problem.

Because $K$ cannot be computed, it can only be approximated from above: a shorter
witnessing program, once found, immediately improves the estimate, but no finite search
can ever certify that no shorter witness exists.

\subsection{The Invariance Theorem}
\label{sec:invariance}

$K$ as defined above depends on the specific reference machine $U$. The Invariance
Theorem states that for any two universal machines $U$ and $V$, there exists a constant
$c_{UV}$, independent of $x$, such that
\[
  K_U(x) \le K_V(x) + c_{UV} \qquad \text{for all } x.
\]
The constant is the length of a program, run on $U$, that simulates $V$: any program of
length $|p|$ that works on $V$ can be run on $U$ by prefixing it with this fixed
simulator, at a fixed additive cost. Consequently, the specific choice of universal
machine -- and hence the specific encoding scheme of Section~\ref{sec:encoding} -- only
ever shifts complexity estimates by a constant that does not grow with the string being
described. This is what licenses treating the particular encoding scheme built for this
project as a legitimate stand-in for "a" universal machine, rather than requiring some
canonical, privileged choice.

\subsection{A Bounded, Computable Approximation}

\texttt{bounded\_kolmogorov} implements a length- and time-bounded estimate of $K(x)$: it
enumerates candidate programs in shortest-first order, decodes and runs each on the UTM
for up to a fixed number of steps, and returns the length of the first one that halts
having emitted exactly the target string. If no such program is found within the length
budget, it returns one more than that budget, signalling that no witness was located
rather than asserting the string's true complexity is that large.

This routine deliberately enumerates \emph{exhaustively} rather than sampling, and
requires halting rather than merely emitting a matching prefix. Both choices follow
directly from what $K(x)$ means: identifying the single shortest program that produces a
given string exactly is a different task from estimating a probability distribution over
which programs are plausible, which is the task the predictor of
Section~\ref{sec:predictor} performs instead. Requiring exact enumeration in shortest-first
order guarantees that the first match found really is shortest; a sampled search could
easily miss the true minimum. The cost of this guarantee is the same combinatorial
explosion that motivated abandoning exhaustive enumeration in the predictor
(Section~\ref{sec:sampling-fails}): \texttt{bounded\_kolmogorov} remains tractable only for
short target strings and small length budgets, and this asymmetry between the two
components -- one exhaustive, one sampled -- is a known, currently unresolved
inconsistency between them, discussed further in Section~\ref{sec:retrospective}.

\section{The Solomonoff Predictor}
\label{sec:predictor}

\subsection{Sequence Prediction and the Universal Prior}

The task is to predict the next bit of an unknown, unbounded binary sequence -- the
"environment" -- having observed some finite prefix $x$ of it so far, with no assumption
about what computable process, if any, is generating that sequence.

Formally, fix a universal prefix machine $U$ and imagine feeding it an infinite stream of
fair coin flips as its program tape. Whatever $U$ outputs under this random input -- a
finite string if $U$ halts, an unbounded one if it does not -- is a single random draw
from the space of all binary sequences, $\{0,1\}^\infty$. This construction induces a
probability measure $M$ on that space: for any finite string $x$, define the
\emph{cylinder set} $\Gamma_x$ as the set of every infinite sequence beginning with $x$,
and let
\[
  M(x) \;:=\; M(\Gamma_x) \;=\; \Pr\big(U\text{'s output, under random coin flips, starts with } x\big).
\]
Because $U$ is deterministic given its input tape, "starts with $x$" is determined
entirely by which prefix of the coin-flip tape $U$ consumes before it has committed to
$x$'s worth of output; different such prefixes are mutually exclusive events, and their
probabilities add. Each candidate prefix $p$ occurs with probability $2^{-|p|}$ (the
probability that $|p|$ fair coin flips come out exactly as $p$), giving the familiar
closed form
\[
  M(x) = \sum_{p \,:\, U(p) \text{ starts with } x} 2^{-|p|}.
\]
The requirement that $U$'s halting programs be \emph{prefix-free} -- no halting program a
prefix of another -- is what makes this sum well defined without double-counting: if $p$
causes $U$ to halt having emitted a prefix of $x$, no proper extension of $p$ can also be
counted, since $U$ never reads past $p$ once it has halted. The same prefix-free property
is precisely Kraft's inequality, $\sum_p 2^{-|p|} \le 1$, restricted to halting programs;
it is what guarantees $M$ is well-behaved as a probability object in the first place, not
what by itself explains why prediction under $M$ converges (Section~\ref{sec:convergence}
addresses that separately). Because some coin-flip tapes drive $U$ into an infinite loop
without ever committing to $n$ full bits of output, $\sum_{x : |x|=n} M(x)$ can be
strictly less than 1; $M$ is therefore a \emph{semimeasure} rather than a full probability
measure, with the missing mass corresponding exactly to non-halting behaviour.

\subsection{From Prior to Prediction}

$M(x)$ is not itself a prediction -- it is the total prior probability mass consistent
with everything observed so far. The predicted probability that the next bit is $1$,
having observed $x$, follows from the ordinary definition of conditional probability,
$P(A\mid B) = P(A\cap B)/P(B)$, applied with $A$ = "the sequence continues with a $1$" and
$B$ = "the sequence starts with $x$":
\[
  P(\text{next} = 1 \mid x) \;=\; \frac{M(x1)}{M(x)},
\]
where $x1$ denotes $x$ with a $1$ appended. Concretely, $M(x)$ is the combined weight of
every candidate program still consistent with the observations so far, and $M(x1)$ is the
weight of the subset of those that additionally agree the next bit is $1$; their ratio is
exactly what the implementation's \texttt{predict\_next} computes.

\subsection{Why Prediction Converges}
\label{sec:convergence}

Suppose the environment is in fact generated by some computable process with a shortest
generating program $p^\ast$ of length $K$. Because $p^\ast$ is, by definition, the process
actually producing the observations, it never contradicts anything seen and is therefore
never eliminated. Its prior weight, $2^{-K}$, is consequently a permanent lower bound on
$M(x)$ for every $x$ the true environment ever produces.

Every other program can be split into those that agree with $p^\ast$ about the next bit at
a given step, and those that do not. A disagreeing program is, by definition, wrong about
something that has not yet been observed; the next contradicting observation removes it
from the pool permanently, and its weight is never recovered. Agreeing programs are never
touched. The total weight belonging to disagreeing programs can therefore only shrink over
time, while the weight belonging to $p^\ast$ and its agreeing companions never does. The
fraction of total surviving weight that could still make an incorrect prediction is
consequently non-increasing, and strictly decreases each time a disagreeing program is
eliminated -- which is precisely the predictor's error, converging toward zero as
observations accumulate. Solomonoff's convergence theorem quantifies the rate of this
process directly in terms of $K$, the complexity of the true generating program: simpler
environments are found, and hence predicted correctly, sooner.

\subsection{Implementation}

The implemented predictor approximates the infinite sum defining $M$ with a finite,
sampled pool of candidate programs (Section~\ref{sec:sampling} describes how this pool is
constructed). Each candidate is stored as a \texttt{\_CandidateProgram} holding its
encoded string, a fixed prior weight, and its output as a list of bits, produced once
during initialisation rather than recomputed on demand.

\paragraph{Precomputing output.} A program's internal computation and the environment's
observation stream operate on unrelated clocks: a program may need an arbitrary, a priori
unknown number of internal steps to emit its next bit, while the environment supplies
exactly one observation per interaction cycle. There is no meaningful way to interleave
the two. Every candidate program is therefore run to completion (or to the step budget)
once, up front, and its emitted output is stored as a fixed list; prediction and update
thereafter only ever index into these precomputed lists, never re-simulate.

\paragraph{Weight normalisation.} The raw prior weight of a program of length $\ell$ is
$2^{-\ell}$; for $\ell$ in the hundreds, this underflows to exactly zero in ordinary
floating-point arithmetic, silently deleting the program's influence. Every consumer of a
weight in this implementation -- \texttt{predict\_next}, \texttt{get\_posterior} -- only
ever computes a \emph{ratio} of weights, never a bare value in isolation. Multiplying every
weight in the pool by the same constant therefore changes none of these ratios. The
implementation exploits this by storing $2^{\,s-\ell}$ instead of $2^{-\ell}$, where $s$ is
the shortest encoding length present in the sampled classes: the shortest program's weight
becomes exactly $1.0$, and every other weight is a representable fraction of it, with no
change to any prediction the system will ever output.

\paragraph{Update and prediction.} On observing bit $b$ at position $n$, every surviving
program whose recorded output at position $n$ disagrees with $b$, or whose output is too
short to have an opinion at position $n$, is discarded, and the total surviving weight is
recomputed as the sum over the remainder. Predicting the following bit sums the weight of
survivors that are both \emph{live} (their output extends past the current position) and
predict a $1$, divided by the total live weight -- programs that have run out of output
are excluded from both numerator and denominator, since including them in the denominator
alone would bias every prediction toward $0$.

\section{Sampling the Hypothesis Space}
\label{sec:sampling}

\subsection{Why Exhaustive Enumeration Fails}
\label{sec:sampling-fails}

An earlier version of the enumerator constructed every machine that fit within a given
length budget directly: every valid combination of alphabet size, blank symbol,
special-state assignment, and transition table. Under the original 7-tuple model this
enumeration was already dominated, by raw count, by the classes with the fewest states,
of order $10^6$ machines within a 40-bit budget on the author's hardware -- and, as
Section~\ref{sec:halt-merge} established, the single largest and cheapest such class was
guaranteed to be computationally inert. Extending the length budget only worsens matters:
the number of possible transition tables for an $(n,k)$-machine grows as
$(4nk)^{nk}$, so even the modest step from a 40-bit to a 115-bit budget under the current
scheme pushes several classes past $10^{30}$ machines, which cannot be enumerated,
let alone simulated, in full.

\subsection{Stratified Sampling by Class}

All machines sharing the same state count $n$ and alphabet size $k$ -- a \emph{class} --
have identical encoded length, and hence identical prior weight. This makes it possible to
draw a program with exactly the correct prior probability without ever enumerating a
class's contents: first choose a class with probability proportional to that class's
\emph{total} prior mass, then choose uniformly among the (astronomically many) machines
within it.

Writing $\text{size}_c$ for the number of machines in class $c$ and $L_c$ for that class's
shared encoding length, the class is chosen with probability
$\text{size}_c \cdot 2^{-L_c} / \text{Total}$, and a machine within it with probability
$1/\text{size}_c$. The probability that this two-stage process selects any particular
machine $p$ in class $c$ is therefore
\[
  P(p) = \frac{\text{size}_c \cdot 2^{-L_c}}{\text{Total}} \times \frac{1}{\text{size}_c}
       = \frac{2^{-L_c}}{\text{Total}},
\]
independent of $\text{size}_c$: it cancels exactly. Every individually drawn machine is
therefore distributed according to the true Solomonoff prior, without the class's
internal contents ever needing to be listed. Drawing $N$ independent samples this way
visits large, low-probability classes rarely and small, high-probability classes often, in
exactly the proportions the theory calls for, at a cost of one class selection and one
uniform draw within it, per sample -- constant work per sample regardless of how large the
underlying class is.

\subsection{Exact Integer Mass Arithmetic}

Class masses combine an astronomically large factor ($\text{size}_c$, which for the
largest classes here exceeds $10^{30}$) with an astronomically small one
($2^{-L_c}$, well below floating-point's smallest representable positive value for large
$L_c$). Computing the product directly in floating point silently produces zero for the
very classes whose size makes them numerically extreme in the first place. Since only
class \emph{ratios} are ever needed (to select a class proportionally, and to choose among
cumulative totals), the implementation instead computes every class's mass as an exact
Python integer, $\text{size}_c \ll (L_{\max} - L_c)$ -- equivalent to
$\text{size}_c \cdot 2^{L_{\max}-L_c}$, a common integer power-of-two rescaling across
every class that cancels in any ratio, computed exactly via a bit shift rather than a
floating-point multiplication. Class selection then uses \texttt{random.randrange} over
the exact integer cumulative total, avoiding the float truncation that a library routine
expecting floating-point weights would otherwise force.

\subsection{Validating the Sampler}

Because a sampler that is subtly biased produces results indistinguishable from a correct
one until specifically checked, the enumerator exposes its own theoretical class-mass
distribution (\texttt{length\_distribution}) and compares it against the lengths actually
observed in a real batch of draws (\texttt{length\_report}). At two million samples with a
115-bit length budget, seventeen distinct classes were reachable; Table~\ref{tab:sampling}
reports a representative subset, comparing observed sampling frequency against the
theoretical expectation computed independently from the class-mass formula.

\begin{table}[h]
\centering
\begin{tabular}{lrrr}
\toprule
Class $(n,k)$ & Encoded length (bits) & Observed frequency & Expected frequency \\
\midrule
$2\times2$ & 25  & 0.521390 & 0.521180 \\
$4\times2$ & 53  & 0.195225 & 0.195442 \\
$2\times4$ & 52  & 0.130020 & 0.130295 \\
$4\times4$ & 112 & 0.048924 & 0.048861 \\
$2\times3$ & 41  & 0.034929 & 0.034785 \\
$3\times2$ & 42  & 0.034848 & 0.034785 \\
$8\times2$ & 115 & 0.014273 & 0.014251 \\
\bottomrule
\end{tabular}
\caption{Observed versus theoretical sampling frequency by class, 2{,}000{,}000 draws,
$L_{\max}=115$. Every observed frequency lies within one standard deviation of its
theoretical expectation, and no length outside the theoretically feasible set was ever
produced.}
\label{tab:sampling}
\end{table}

The three cheapest classes -- $2\times2$, $4\times2$, and $2\times4$ -- together account
for $84.7\%$ of all prior mass at this length budget, a direct, measured illustration of
how sharply the Solomonoff prior favours short descriptions.

\section{Empirical Results}
\label{sec:results}

\subsection{Convergence on Structured Sequences}

Fed the periodic sequence $0,1,0,1,0,1,\ldots$, the predictor's estimate of
$P(\text{next}=1)$ evolved as $0.500 \to \sim\!0 \to 0.518 \to 0.769 \to 0.137 \to 1.000
\to \sim\!0 \to 1.000$ over the first eight observations: an initial guess of pure
uncertainty, a confident (and correct) guess of "more zeros" after the first zero is seen,
a period of readjustment as the alternating pattern reveals itself, and a locked-on,
correct prediction from the sixth observation onward. The corresponding survivor count
falls from roughly one million candidate programs to a stable, non-collapsing few over the
same span -- the signature the theory of Section~\ref{sec:convergence} predicts: mass
concentrating on the (many, weight-equivalent, all period-2-consistent) surviving
hypotheses rather than being exhausted entirely.

\subsection{Failure on Sequences Requiring Counters}

Not every computable sequence is found. A hand-constructed "doubling blocks" sequence
(runs of $0$s and $1$s of length $1,1,2,2,4,4,8,8,\ldots$) drove the survivor count to
exactly zero after roughly fifteen observations, at which point \texttt{predict\_next}
falls back to an uninformative $0.5$. This is a distinct failure mode from a wrong but
confident prediction: it indicates that no program in the two-million-sample pool
generates this sequence, not that the predictor converged incorrectly. The reason is
structural rather than a matter of sample size alone -- generating growing run-lengths
requires a program that can count, which in turn requires enough working states and
transitions to implement a counter, and the classes capable of this are precisely the
ones the prior assigns the least mass to. This is an honest, expected limitation of an
$l,t$-bounded approximation to Solomonoff induction: the theory guarantees convergence
whenever the true generating program lies within the sampled hypothesis space, and offers
no guarantee otherwise.

\subsection{The Limits of a Deterministic Program Model}

A genuinely stochastic source, such as a biased coin with $P(1)=0.7$, cannot be modelled
by this implementation at all, for a structural reason rather than an insufficient sample
budget: every candidate program runs on an empty input tape and produces one fixed,
deterministic output string. There is no true generating \emph{program} for a stochastic
source among deterministic candidates -- observing enough bits from a Bernoulli process
eliminates every fixed-output hypothesis in turn, and the survivor pool collapses to
nothing regardless of how it is sampled. The theoretically correct extension, present in
Hutter's original formulation but not implemented here, equips the universal machine with
a second, read-only random input tape that programs may consult during execution; $M$
then becomes a probability-weighted sum over program \emph{behaviours} rather than a hard
consistency filter over program \emph{outputs}, and a single program can legitimately
represent a stochastic process by reading random bits and using them to bias its emitted
output. Implementing this extension -- a second input tape, multiple simulation runs per
candidate with different random seeds, and a corresponding change to how a program's
weight contributes to $M(x)$ -- is identified here as concrete future work rather than
attempted in the current implementation.


\end{document}